\section{An exact low-height exclusion from the original theta boundary}
\label{sec:quartet-low-height-theta}

This calculation keeps the original boundary $f_0=\Theta\phi_*$,
its Mellin transform $g=2\xi$, and its unscaled positive kernel
$\psi$.  In particular, the constant $g(1)=1$ below is a value
calculated from that boundary, not a normalization imposed on $g$.
Every estimate is an inequality of actual integrals.  The numerical
lower bound at the end is a rational number proved positive by exact
arithmetic.

\subsection{The original source and its exact endpoint mass}

Use the functions and coordinates in (HCA.1)--(HCA.2):
\[
\begin{gathered}
 \vartheta(x)=\sum_{n\in\mathbb Z}e^{-\pi n^2x^2}
       =1+2\sum_{n\geq1}e^{-\pi n^2x^2},\\
 \phi_*(x)=(4\pi^2x^4-6\pi x^2)e^{-\pi x^2},\\
 f_0(x)=\Theta\phi_*(x)=2\sum_{n\geq1}\phi_*(nx),\\
 \psi(y)=e^{y/2}f_0(e^y),\\
 g(s)=\mathcal M f_0(s)=\int_0^\infty f_0(x)x^s\frac{dx}{x}
     =\int_{\mathbb{R}}e^{(s-1/2)y}\psi(y)\,dy.
\end{gathered}
\tag{LHT.1}
\]
Here $x>0$ and $y=\log x\in\mathbb{R}$.  The original theta identity
is $\vartheta(x)=x^{-1}\vartheta(x^{-1})$.  This is the same Gaussian
theta identity used in the source proof of
$\psi(-y)=\psi(y)$ in (HCA.1); no different theta function is
introduced here.  Its series and every differentiated series converge
uniformly on compact subsets of $x>0$.  Differentiating them gives
\[
 \begin{split}
 \vartheta'(x)&=-4\pi x\sum_{n\geq1}n^2e^{-\pi n^2x^2},\\
 (x^2\vartheta'(x))'
 &=8\pi^2x^4\sum_{n\geq1}n^4e^{-\pi n^2x^2}
       -12\pi x^2\sum_{n\geq1}n^2e^{-\pi n^2x^2}
 =f_0(x).
 \end{split}
\tag{LHT.2}
\]
For $x\geq1$, every term of $f_0(x)$ is positive, since
$4\pi^2n^4x^4-6\pi n^2x^2
=2\pi n^2x^2(2\pi n^2x^2-3)>0$.  The evenness of $\psi$ therefore
gives $\psi(y)>0$ for every real $y$.  On $y\geq0$ we retain the
termwise upper bound
\[
 0<\psi(y)\leq
 8\pi^2\sum_{n\geq1}n^4
       \exp\bigl(\tfrac92y-\pi n^2e^{2y}\bigr).
\tag{LHT.3}
\]
Together with evenness, this is an integrable majorant after
multiplication by any fixed polynomial in $y$ and any fixed
$e^{b|y|}$.  For example, for $y\geq0$ the sum is bounded by
$\bigl(\sum_{n\geq1}n^4e^{-\pi(n^2-1)}\bigr)
e^{9y/2-\pi e^{2y}}$ times $8\pi^2$.
All Mellin integrals and differentiations below consequently converge
absolutely; termwise upper bounds on their nonnegative integrands
also follow directly from monotone convergence.

The limits in (LHT.2) can be evaluated before integrating.
At $x\to\infty$, the differentiated theta series gives
$x^2\vartheta'(x)\to0$.  Differentiating the theta identity itself
gives the exact equation
\[
 x^2\vartheta'(x)
       =-\vartheta(x^{-1})-x^{-1}\vartheta'(x^{-1}).
\tag{LHT.4}
\]
As $x\to0^+$ the two terms on the right tend to $-1$ and $0$,
respectively.  Integrating (LHT.2) first on a compact interval and
then taking its endpoints to $0$ and $\infty$ is justified by the
absolute convergence just proved.  It follows that
\[
 g(1)=\int_0^\infty f_0(x)\,dx
     =\bigl[x^2\vartheta'(x)\bigr]_{0^+}^{\infty}=1.
\tag{LHT.5}
\]
This calculation does not interchange the integrals of the signed
individual theta terms on the whole positive axis.

Define two actual, unscaled source moments
\[
 \begin{split}
 M_0&=\int_{\mathbb{R}}\psi(y)\,dy=g(1/2),\\
 G_2&=\int_{\mathbb{R}}y^2e^{y/2}\psi(y)\,dy
       =2\int_0^\infty y^2\cosh(y/2)\psi(y)\,dy=g''(1).
 \end{split}
\tag{LHT.6}
\]
Both numbers are strictly positive.  For real $z$,
$\cosh z-1\leq \frac12 z^2\cosh z$: for $z\geq0$ this follows from
$\cosh z-1=\int_0^z(z-v)\cosh v\,dv
\leq\cosh z\int_0^z(z-v)\,dv$, and both sides are even.
Equations (LHT.5)--(LHT.6) therefore imply
\[
 1-M_0
 =2\int_0^\infty (\cosh(y/2)-1)\psi(y)\,dy
 \leq \frac18 G_2,
 \qquad M_0\geq1-\frac18G_2.
\tag{LHT.7}
\]

\subsection{A rational upper bound on the unscaled second moment}

We now bound $G_2$ without a floating-point computation or a
zero-location theorem.  For $y\geq0$ use
$\cosh(y/2)\leq e^{y/2}$ in (LHT.3), and then put $x=e^y$.
The resulting nonnegative sum and integrals give
\[
 \begin{split}
 G_2
 &\leq16\pi^2\sum_{n\geq1}n^4
      \int_0^\infty y^2e^{5y-\pi n^2e^{2y}}\,dy\\
 &=16\pi^2\sum_{n\geq1}n^4
      \int_1^\infty (\log x)^2x^4e^{-\pi n^2x^2}\,dx.
 \end{split}
\tag{LHT.8}
\]
For $x\geq1$, $\log x\leq x-1$.  Set $v=x-1\geq0$ and use
$(1+v)^2\geq1+2v$.  Expanding the remaining polynomial, with no
change to the coordinate or the original moment, yields
\[
 \begin{split}
 G_2
 &\leq16\pi^2\sum_{n\geq1}n^4e^{-\pi n^2}
         \int_0^\infty v^2(1+v)^4e^{-2\pi n^2v}\,dv\\
 &=\sum_{n\geq1}e^{-\pi n^2}
     \sum_{r=0}^{4}
       \frac{16\binom4r(r+2)!}
          {2^{r+3}\pi^{r+1}n^{2r+2}}.
 \end{split}
\tag{LHT.9}
\]
The integral identity used here is
$\int_0^\infty v^j e^{-av}\,dv=j!a^{-j-1}$ for $a>0$ and
nonnegative integer $j$; it follows by $j$ integrations by parts,
with base integral $1/a$ and zero endpoint terms.

The elementary inequality $\pi>3$ and $n\geq1$ bound the five
terms of the inner sum, respectively, by
\[
 \frac43,\qquad \frac83,\qquad \frac83,
 \qquad \frac{40}{27},\qquad \frac{10}{27};
 \qquad
 \frac43+\frac83+\frac83+\frac{40}{27}+\frac{10}{27}
       =\frac{230}{27}.
\tag{LHT.10}
\]
For the exponential bound, the exact finite Taylor sum gives
\[
 e^3>\sum_{j=0}^{8}\frac{3^j}{j!}
      =\frac{89641}{4480}>\frac{89600}{4480}=20.
\tag{LHT.11}
\]
Thus $e^{-\pi n^2}<20^{-n^2}\leq20^{-n}$ for $n\geq1$, and
\[
 \sum_{n\geq1}e^{-\pi n^2}<
 \sum_{n\geq1}20^{-n}=\frac1{19}.
 \tag{LHT.12}
\]
Combining (LHT.9)--(LHT.12) proves the entirely rational bound
\[
 0<G_2=g''(1)<\frac{230}{27\cdot19}
                  =\frac{230}{513}.
\tag{LHT.13}
\]
In particular this proof includes every $n\geq1$ and the entire
positive integration axis.  There is no unestimated theta tail.

\subsection{Uniform exclusion in the closed low-height rectangle}

Write $s=\frac12+a+it$, where $a,t\in\mathbb{R}$.
Evenness and reality of $\psi$ in (LHT.1) give
\[
 \operatorname{Re}g(\tfrac12+a+it)
 =2\int_0^\infty\psi(y)\cosh(ay)\cos(ty)\,dy.
\tag{LHT.14}
\]
The real inequality $\cos z\geq1-z^2/2$ follows from
$1-\cos z=2\sin^2(z/2)\leq z^2/2$.  If $|a|\leq1/2$, then
$1\leq\cosh(ay)\leq\cosh(y/2)$ for $y\geq0$.
Since every factor $\psi(y)\cosh(ay)$ is nonnegative, substitution
of this cosine inequality in (LHT.14) gives
\[
 \begin{split}
 \operatorname{Re}g(\tfrac12+a+it)
 &\geq 2\int_0^\infty\psi(y)\cosh(ay)\,dy
       -t^2\int_0^\infty y^2\psi(y)\cosh(ay)\,dy\\
 &\geq M_0-\frac{t^2}{2}G_2.
 \end{split}
\tag{LHT.15}
\]
For $|t|\leq2$, equations (LHT.7) and (LHT.13) now imply
\[
 \boxed{\quad
 \operatorname{Re}g(s)
 \geq1-\frac{17}{8}G_2
 >1-\frac{17}{8}\frac{230}{513}
 =\frac{97}{2052}>0,
 \qquad 0\leq\operatorname{Re}s\leq1,
 \quad |\operatorname{Im}s|\leq2.
 \quad}
\tag{LHT.16}
\]
Consequently $g(s)=2\xi(s)$ has no zero in that closed rectangle.
This is proved directly from the actual theta source, without
assuming the Riemann hypothesis or using numerical verification
of its low zeros.

The resulting map on the programme's original local packet objects
is also explicit.  At every point $s_0$ of this rectangle,
$g(s_0)\ne0$, so multiplication by $g$ on the local holomorphic
ring $\mathcal O_{\mathbb{C},s_0}$ is an isomorphism with inverse
multiplication by the holomorphic germ $1/g$.  On the rectangle
itself the preceding calculation proves
\[
 |1/g(s)|\leq\frac1{\operatorname{Re}g(s)}
                  <\frac{2052}{97}.
\tag{LHT.17}
\]
It follows that
$\mathcal O_{\mathbb{C},s_0}/(g)=0$ at every such $s_0$.
Under the programme's already constructed Mellin--jet isomorphism
to the actual packet summands, each of these local summands therefore
vanishes.  No proper-source kernel is quotiented out by
this argument.  In particular an actual upper-half-plane offcritical
centre obeys
\[
 g(\rho)=0,\quad
 \rho=\frac12+\delta+i\gamma,\quad
 0<\delta<\frac12,\quad \gamma>0
 \quad\Longrightarrow\quad \gamma>2.
\tag{LHT.18}
\]
Its reflected and conjugate centres retain their original orders,
jets, and Taylor units.
